\chapter{Metodologia}

\section{Configuração do Ambiente}

Os dados apresentados em \ref{chap:resultados} foram gerados a partir de uma
máquina com as seguintes configurações:

\begin{table}[ht]
    \centering
    \caption{Especificações Técnicas do Ambiente de Testes}
    \label{tab:especificacoes_sistema}
    \begin{tabularx}{\textwidth}{l X}
        \toprule
        \textbf{Componente} & \textbf{Descrição} \\
        \midrule
        \textbf{Hardware} & \\
        Modelo & Lenovo IdeaPad 320-15IKB (Type 80YH) \\
        Processador & Intel Core i7-7500U (2 núcleos, 4 threads) \\
        Clock & 2,70 GHz (Base) / 3,50 GHz (Turbo Max) \\
        Memória RAM & 16 GiB \\
        Armazenamento & HDD Seagate 2,0 TB (5400 RPM, SATA 6,0 Gb/s) \\
        GPU Integrada & Intel HD Graphics 620 \\
        GPU Dedicada & NVIDIA GeForce 940MX (Arquitetura Maxwell) \\
        \midrule
        \textbf{Software} & \\
        Sistema Operacional & Linux Mint 22.1 "Xia" (64 bits) \\
        Kernel & 6.8.0-87-generic (x86\_64) \\
        Base de Distro & Ubuntu 24.04 LTS (Noble Numbat) \\
        Ambiente Gráfico & Cinnamon 6.4.6 (Servidor X.Org 21.1.11) \\
        JVM de Execução & OpenJDK 21 \\
        Compilador & GCC v13.3.0 \\
        Drivers de Vídeo & i915 (Intel) / Nouveau (NVIDIA) / Mesa v24.0.9 \\
        \bottomrule
    \end{tabularx}
    \vspace{2pt}
    \small
    \centering
\end{table}

Para garantir maior confiabilidade e precisão dos dados, a literatura recomenda
desabilitar o Intel Turbo Boost e o \textit{Hyper-Threading}
\cite{vms-hot-and-cold, steady-state-assessment}. Essas funcionalidades podem
introduzir ruído nas medições, pois alteram dinamicamente o comportamento da
CPU ao longo das execuções. Por esse motivo, todos os \textit{benchmarks}
foram executados com esses recursos desativados. Além disso, cada execução foi
fixada a um único \textit{core}. Embora alguns \textit{benchmarks} possam se
beneficiar de paralelismo, essa decisão reduziria a comparabilidade entre os
resultados e adicionaria \textit{overhead} energético decorrente da própria
administração das \textit{threads}.

\section{Medição de Energia}

Para medir o consumo de energia, há duas metodologias predominantes na
literatura. A primeira consiste no uso de \textit{hardware} suplementar para
medir o consumo diretamente da tomada, resultando em uma medida fiel do consumo
real de energia \cite{energy-java-collections-hasan}. A segunda abordagem
consiste em utilizar estimativas de consumo geradas pela própria máquina. Neste
trabalho, optou-se pela segunda metodologia. A ferramenta utilizada para
estimar o consumo de energia durante a execução dos \textit{benchmarks} foi o
Intel Running Average Power Limit (RAPL) \cite{rapl}. O RAPL é uma
funcionalidade introduzida nas CPUs da Intel a partir da arquitetura Sandy
Bridge. Ele utiliza registradores de 32 \textit{bits} para armazenar a
informação do consumo de energia desde a inicialização da máquina; esses
registradores são atualizados aproximadamente a cada 1 milissegundo. Por meio
dessa ferramenta, é possível obter estimativas do consumo de energia de
diferentes componentes, como o pacote completo da CPU, seus \textit{cores} e a
DRAM. O RAPL apresenta precisão suficiente para servir como uma aproximação
fiel do consumo real de energia: \citeonline{rapl-in-action} encontraram
correlação de 0,99 entre o valor estimado pelo RAPL e o valor real medido
externamente. Além disso, os autores mostram que a sobrecarga introduzida pelo
próprio mecanismo é negligenciável. Outra vantagem do uso do RAPL em relação a
métodos diretos é sua maior granularidade, permitindo focalizar partes
específicas do \textit{hardware}, enquanto medições externas normalmente
observam apenas o sistema como um todo. Sendo assim, o RAPL constitui um método
confiável e conveniente para obter medições de energia.

É possível acessar os registradores do RAPL diretamente para obter os valores
de energia. Entretanto, em vez desse método, foi empregada a ferramenta JRAPL
\cite{jrapl}. O JRAPL é um \textit{wrapper} que permite acessar os
registradores diretamente a partir de código Java. A sobrecarga decorrente
desse acesso é negligenciável, de modo que seu uso não compromete as medições
\cite{programming_lang_energy}. A integração direta com o código Java é uma
vantagem importante do JRAPL em relação ao acesso manual aos registradores, pois
elimina problemas de sincronização entre o código dos \textit{benchmarks} e a
coleta das estimativas, garantindo que as leituras sejam obtidas imediatamente
antes e imediatamente depois das execuções.

\section{Execução de benchmarks na JVM}

A execução de \textit{benchmarks} em Java apresenta desafios adicionais em
comparação a outras linguagens. Isso se deve principalmente ao ambiente de
execução da linguagem, isto é, à Java Virtual Machine (JVM).

Ao contrário de linguagens compiladas estaticamente, o comportamento de uma
aplicação Java pode se alterar ao longo da própria execução. Esse efeito
decorre da compilação \textit{just-in-time} (JIT), do perfilamento adaptativo
do código, do carregamento dinâmico de classes e da atuação do coletor de lixo.
Como consequência, as primeiras iterações de um \textit{benchmark}
frequentemente representam uma fase de aquecimento, e não o comportamento
estável do sistema \cite{vms-hot-and-cold, steady-state-assessment}. Além
disso, mesmo em ambientes fortemente controlados, nem toda combinação entre
\textit{benchmark} e JVM necessariamente atinge um regime estável de
desempenho, o que torna arriscado assumir, \textit{a priori}, que poucas
iterações sejam suficientes para caracterizar o programa
\cite{vms-hot-and-cold}.

Com base nessas observações, a metodologia adotada neste trabalho separa
explicitamente as iterações de aquecimento das iterações utilizadas na análise
final. Para isso, foi empregada uma versão estendida do DaCapo Chopin
\cite{dacapochopin}, instrumentada com um \textit{callback} próprio chamado
\texttt{EnergyCallback}. Esse \textit{callback} coleta as leituras de energia
imediatamente antes e imediatamente depois de cada iteração, utilizando o
JRAPL, e registra também o tempo decorrido informado pelo próprio
\textit{harness}. Dessa forma, cada observação reúne, de maneira sincronizada,
o tempo de execução e a energia estimada nos domínios DRAM, CPU e pacote do
processador.

\section{Configurações do DaCapo}

O DaCapo Chopin foi escolhido por disponibilizar um conjunto diverso de cargas
de trabalho Java, o que reduz o risco de se tirar conclusões a partir de um
conjunto excessivamente restrito de aplicações \cite{dacapochopin}. Embora a
suíte completa contenha 22 \textit{benchmarks}, neste trabalho foram
considerados os 16 \textit{benchmarks} que puderam ser construídos e executados
corretamente no ambiente experimental. Todas as execuções utilizaram o perfil
de carga \texttt{small}, mantendo a mesma classe de entrada ao longo de todos
os experimentos para preservar a comparabilidade entre as medições.

As execuções foram automatizadas por dois scripts. O script
\texttt{run\_energy.sh} realiza uma execução individual do \textit{benchmark}
instrumentado, enquanto o script \texttt{run\_sweep.sh} executa, de forma
sequencial, todas as combinações de parâmetros definidas para o experimento. Em
cada execução, o processo Java é fixado ao núcleo 0 por afinidade de CPU, o
\textit{Hyper-Threading} é desabilitado, o Intel Turbo Boost é desligado e, ao
final, as configurações originais da máquina são restauradas. Esse controle
busca reduzir fontes externas de variabilidade, em consonância com a literatura
sobre avaliação de desempenho em JVMs
\cite{vms-hot-and-cold, steady-state-assessment}.

Além do isolamento do processo, foram variados dois fatores experimentais
principais: frequência do processador e tamanho de \textit{heap}. A frequência
foi configurada manualmente, por meio do \texttt{cpupower}, nos valores
400 MHz, 800 MHz, 1{,}2 GHz, 1{,}6 GHz, 2{,}0 GHz, 2{,}4 GHz e 2{,}7 GHz. A
escolha por varrer a frequência decorre do interesse em observar a interação
entre o comportamento da aplicação e mecanismos de escalonamento dinâmico de
frequência, aspecto já explorado em trabalhos baseados em JRAPL \cite{jrapl}.
Para cada \textit{benchmark}, foi definido um conjunto próprio de tamanhos de
\textit{heap}, sempre com \texttt{-Xms} e \texttt{-Xmx} configurados com o
mesmo valor, evitando redimensionamentos dinâmicos de memória ao longo da
execução e tornando mais controlado o espaço de memória disponível para a JVM.

Nas execuções consideradas neste relatório, utilizou-se o coletor de lixo G1 e
o binário do OpenJDK 21. Para cada combinação de \textit{benchmark}, frequência
e \textit{heap}, foram realizadas 10 repetições independentes. Em cada
repetição, o DaCapo executa iterações de aquecimento e uma iteração medida, com
limite máximo de 40 iterações por tentativa. Quando o próprio \textit{harness}
informa falha de convergência, a repetição é descartada e executada novamente,
de modo que apenas repetições concluídas com sucesso sejam mantidas no conjunto
de dados. Esse procedimento reduz a chance de incorporar à análise execuções
ainda instáveis ou representativas apenas de uma fase transitória da JVM.

\section{Tratamento e análise dos dados}

O \textit{callback} implementado neste trabalho grava, ao final de cada
iteração, uma linha em um arquivo CSV contendo o nome do \textit{benchmark}, a
informação de validade, a indicação de aquecimento, o tempo decorrido em
milissegundos, o soquete medido, o tamanho de \textit{heap}, o coletor de
lixo, a frequência do processador e as energias estimadas para DRAM, CPU e
pacote. Como a máquina experimental possui um único soquete, cada iteração
medida corresponde, na prática, a uma única linha de dados no arquivo
\texttt{energy.csv}. Entre os domínios coletados, a energia do pacote foi
adotada como métrica principal de comparação, pois representa o consumo
agregado do processador durante a execução e foi a base para o cálculo das
métricas derivadas empregadas na análise.

O pós-processamento foi realizado em Python. O script
\texttt{process\_energy.py} primeiro descarta todas as linhas marcadas como
\textit{warmup}, preservando apenas as iterações efetivamente utilizadas para
comparação. Em seguida, são calculadas métricas derivadas a partir da energia
do pacote e do tempo de execução: o \textit{Energy-Delay Product} (EDP),
definido como \(E \times T\), e o \(EDP^2\), definido como \(E^2 \times T\).
Essas métricas são úteis porque combinam energia e desempenho em uma mesma
medida, permitindo observar situações em que reduzir o tempo de execução não
implica, necessariamente, reduzir o consumo energético, como já apontado por
\citeonline{programming_lang_energy}.

Após esse passo, os dados são agrupados por \textit{benchmark}, tamanho de
\textit{heap} e frequência do processador. Para cada grupo, o script calcula a
média e o intervalo de confiança de 95\% por meio de \textit{bootstrap} com
1000 reamostragens. Essa escolha permite resumir os resultados sem assumir
previamente uma distribuição paramétrica específica para as medições. Por fim,
o script \texttt{generate\_graphs.py} normaliza os tamanhos de \textit{heap}
em múltiplos do menor valor avaliado para cada \textit{benchmark} e gera
gráficos de energia, tempo, EDP e \(EDP^2\) em função da frequência do
processador, incluindo bandas de confiança. Esses gráficos constituem a base da
análise apresentada no Capítulo \ref{chap:resultados}.
